%%%%%%%%%%%%%%%%%%%%%%%%%%%%%%%%%%%%%%%%%%%%%%
% Journal of Law and Courts target manuscript.
% Official Cambridge/Overleaf small-template source uses cup-journal.
% JLC author instructions: anonymous review, Chicago author-date citations,
% tables/figures near first reference, and replication materials at acceptance.
%%%%%%%%%%%%%%%%%%%%%%%%%%%%%%%%%%%%%%%%%%%%%%
\documentclass[
  journal=small,
  manuscript=research-article,
  year=2026
]{cup-journal}

\usepackage{amsmath}
\usepackage[nopatch]{microtype}
\usepackage{booktabs}
\usepackage{graphicx}
\usepackage{longtable}
\usepackage[authoryear,round]{natbib}
\usepackage{url}
\usepackage{xcolor}
\usepackage[hidelinks]{hyperref}

\makeatletter
\newenvironment{inlinefigureblock}{%
  \par\smallskip
  \begingroup
  \centering
}{%
  \par\smallskip
  \endgroup
}

\newcommand{\inlinefigurecaption}[2]{%
  \refstepcounter{figure}%
  \label{#2}%
  \vskip\abovecaptionskip
  \@makecaption{\figurename~\thefigure}{#1}%
  \vskip\belowcaptionskip
}
\makeatother

\title{Emergency Review and Democratic Constitutionalism: A Comparative Simulation Framework}

\author{Anonymous Author}
\affiliation{Affiliation withheld for review}
\date{}

\keywords{constitutional review, judicial politics, comparative courts, institutional design, simulation}

\begin{document}
\maketitle

\begin{abstract}
Emergency constitutional review is where many court-reform debates become most concrete: interim relief can protect rights before merits review, but it can also produce opaque, merits-displacing vetoes. This article uses a comparative simulation framework to study emergency-review design, with appointment systems, councils, overrides, weak-form review, suspended declarations, pre-enactment review, abstract review, ombudsman access, constitutional public defenders, rights-impact statements, and mandatory legislative response cycles included as comparison infrastructure rather than equal empirical targets. The empirical contribution is limited to source-specific diagnostic checks with documented sources, denominators, and direct target analogues; broader comparative rows remain transparent stress-test assumptions. The synthetic finding is narrower: rules requiring written emergency reasons, merits follow-through, or no relief without merits connection reduce modeled emergency-process irregularity more directly than broad appointment or size reforms, but the relevant evaluation depends on rights-claimant success, compliance, public-legitimacy proxies, implementation cost, legal-transplant feasibility, political-culture sensitivity, and veto-relocation risk. The design recommendations are therefore conditional. Emergency review improves democratic constitutionalism only when it protects urgent rights without hiding final policy control in opaque or weakly accountable institutional stages.
\end{abstract}

\section{Introduction}

Constitutional review systems face a recurring design problem: courts need enough independence to enforce constitutional limits, but enough accountability and procedural discipline to avoid becoming unreviewable political actors. The classic countermajoritarian framing emphasizes the difficulty of judicial invalidation in a democracy \citep{bickel1962}, and court-skeptical accounts press the democratic case for limiting judicial supremacy \citep{tushnet1999}. Comparative constitutional scholarship shows that review can be organized through many institutional forms rather than a single supreme-court template \citep{stoneSweet2000,ginsburg2003}. Law-and-courts scholarship also treats courts as institutions whose behavior depends on selection rules, legal hierarchy, compliance expectations, and political context \citep{epsteinKnight1998,helmke2005,vanberg2005,staton2010}.

The premise of this article is that those institutional choices should be analyzed jointly, but the paper's substantive focus is narrower than the full simulator. Emergency review provides the clearest test case because it combines urgent rights protection, interim policy control, procedural opacity, and weak merits follow-through in the same institutional setting. Appointment method affects ideological composition, but its effects depend on tenure, removal standards, court size, docket control, voting thresholds, and emergency procedure. A strict supermajority rule may improve stability in ordinary merits cases but weaken claimant success when the docket contains high-risk legislation. A constitutional council may reduce ex post invalidation but shift costs into upstream screening. A reason-giving emergency procedure may reduce emergency-process irregularity while slowing response time. These are system-level interactions.

The contribution is a computational simulation framework for comparing emergency-review and constitutional-review design bundles. The model is deliberately stylized. It does not forecast the U.S. Supreme Court, the German Federal Constitutional Court, or any other real tribunal. Instead, it provides a structured way to ask how design bundles behave under shared dockets and shared worlds. Institution-specific empirical models can represent one court in greater detail; this article asks which emergency-review reforms and adjacent institutional architectures are worth comparing before that institution-specific work begins.

The framework can also use legislative-output profiles as subject-matter-preserving docket stress signals. That bridge makes it possible to compare review designs against the kinds of laws different legislative systems produce without treating legislative and judicial institutions as the same model. The central substantive question is modest but important: when does constitutional review improve democratic constitutionalism rather than merely relocating veto power? Appointment method, court size, and term length matter, but so do front-end rights review, weak-form declarations, abstract review, emergency procedure, lower-court routing, voting thresholds, implementation cost, and compliance culture. The article therefore treats simulation as a comparative institutional audit, not as a predictive model of existing courts.

The strongest claim tested below is therefore not that a simulator can identify the best constitutional court. It is that emergency-process reforms can be compared against appointment, council, override, and weak-form reforms in a shared-world design, and that the most informative outputs are second-order tradeoffs rather than the obvious fact that a transparency rule improves a transparency index. The rest of the article is organized for a law-and-courts audience. Section~\ref{sec:infrastructure} defines constitutional review as institutional infrastructure. Section~\ref{sec:framework} describes the simulation framework. Section~\ref{sec:calibration} explains calibration, uncertainty, and replication. Section~\ref{sec:experiments} summarizes the comparative experiments. Section~\ref{sec:implications} identifies implications for institutional design. Section~\ref{sec:limits} discusses limitations and next steps.

\section{Constitutional Review as Institutional Infrastructure}
\label{sec:infrastructure}

A constitutional review system is more than an apex court. It includes appointment and replacement rules, lower-court pathways, docket filters, emergency relief procedures, merits review, voting thresholds, opinion coalitions, recusal practices, legislative override rules, and enforcement channels. Treating these features as separate design variables makes sense analytically, but reform proposals usually bundle them. The simulation therefore models review as an infrastructure problem: a constitutional system routes disputes through a sequence of institutions, and the consequences depend on how those routing rules interact.

Table~\ref{tab:design-space} summarizes the current design space. The categories are intentionally broad enough to cover U.S.-style discretionary appellate review, specialized constitutional complaints, pre-enactment councils, mixed abstract and concrete review, declaration-oriented parliamentary review, and supranational treaty courts.

\begin{longtable}{p{0.32\linewidth}p{0.58\linewidth}}
\caption{Modeled constitutional-review design space.}
\label{tab:design-space}\\
\toprule
Design family & Modeled variants \\
\midrule
\endfirsthead
\toprule
Design family & Modeled variants \\
\midrule
\endhead
Appointment method & Presidential/senate appointment, supermajority commission, merit sortition, legislative selection, and popular-retention variants. \\
Court size & 9-, 11-, 12-, 15-, 16-, 18-, 27-, 34-, and 46-member designs for experimental and real-world benchmark scenarios. \\
Terms and removal & Life tenure, nonrenewable long terms, short renewable terms, impeachment-only removal, misconduct commission removal, and retention-election accountability. \\
Recusal & Self-policing, mandatory conflict recusal, random substitution, and strict transparency. \\
Docket control & Discretionary certiorari, leave to appeal, complaint admissibility, referral-gated review, and mandatory review with filters. \\
Emergency docket & Fast shadow docket, reasoned emergency panels, full-court emergency review, and merits follow-up. \\
Doctrine & Speech, equality, criminal procedure, federalism, election law, emergency powers, and administrative state claims. \\
Policy domains & Civil rights, speech/religion, criminal justice, federalism, governance, elections, emergency/security, economic regulation, and administration. \\
Jurisdiction and lower courts & Federal, state, and mixed state-federal disputes; district-only, circuit-panel, en banc, state high-court, and state-federal split paths. \\
Voting threshold & Simple majority, supermajority invalidation, concurrent-majority logic, and high constitutional threshold. \\
Review structure & Full court, panel plus en banc review, dual supreme courts, cross-checking courts, and constitutional council review. \\
Legislative override & No override, supermajority override, delayed majority override, and referendum override. \\
Real-world benchmark presets & U.S. Supreme Court, German Federal Constitutional Court, French Constitutional Council, Supreme Court of Canada, South African Constitutional Court, UK Supreme Court, India Supreme Court, Brazil STF, ECHR, and CJEU. \\
Synthetic mechanism tests & Weak-form review, suspended declarations, explicit override clauses, pre-enactment review, abstract review, ombudsman-triggered review, constitutional public defenders, rights-impact statements, and mandatory legislative response cycles. \\
Feasibility modifiers & Legal-transplant feasibility, implementation cost, and political-culture sensitivity are modeled as diagnostic constraints rather than empirical country estimates. \\
\bottomrule
\end{longtable}

The central design question is not whether one feature is good or bad in isolation. It is whether a bundle produces an acceptable balance among legal stability, rights protection, democratic responsiveness, legitimacy, conflict management, administrative burden, and enforceability.

The project therefore separates three evidentiary categories. Real-world presets encode observable institutional rules and cost/intake context for named systems, but they are not fitted replicas. Synthetic mechanisms isolate reform ideas that may have comparative analogues but are not country presets. Speculative design recommendations are conditional inferences from those synthetic comparisons; they are not empirical claims about the systems that inspired a mechanism.

\section{Simulation Framework}
\label{sec:framework}

The simulator uses named scenarios, randomized worlds, shared dockets, campaign sweeps, directional metrics, structured reports, and provenance records. Each case has a generated policy position, doctrine area, policy domain, jurisdiction, lower-court path, rights-threat signal, public support, legislative mandate, urgency, legal ambiguity, constitutional salience, lower-court conflict, lower-court panel skew, state-federal tension, intercourt conflict, certiorari pressure, time-to-review, executive pressure, conflict-of-interest risk, public-trust context, litigant capacity, public-interest support, and government repeat-player advantage.

Figure~\ref{fig:model-pipeline} summarizes the model pipeline. Legislative-output profiles enter only as structured stress signals that alter the generated constitutional-review docket.

\begin{figure}[hbt!]
\centering
% Auto-generated by paper/scripts/generate_figures.py
\begingroup
\setlength{\unitlength}{1mm}
\begin{picture}(128,66)
\scriptsize
\put(2,45){\framebox(22,12){\shortstack[c]{Scenario\\catalog}}}
\put(28,45){\framebox(22,12){\shortstack[c]{World and\\case generation}}}
\put(54,45){\framebox(22,12){\shortstack[c]{Institutional\\routing}}}
\put(80,45){\framebox(22,12){\shortstack[c]{Decision\\process}}}
\put(106,45){\framebox(20,12){\shortstack[c]{Diagnostics\\and reports}}}
\put(24,51){\vector(1,0){4}}
\put(50,51){\vector(1,0){4}}
\put(76,51){\vector(1,0){4}}
\put(102,51){\vector(1,0){4}}
\put(28,26){\framebox(22,10){\shortstack[c]{Legislative\\outputs}}}
\put(39,36){\vector(0,1){9}}
\put(54,28){\framebox(22,10){\shortstack[c]{Intake and\\access filters}}}
\put(80,28){\framebox(22,10){\shortstack[c]{Emergency\\and merits}}}
\put(106,28){\framebox(20,10){\shortstack[c]{Compliance\\reaction cost}}}
\put(50,33){\vector(1,0){4}}
\put(76,33){\vector(1,0){4}}
\put(102,33){\vector(1,0){4}}
\put(116,28){\vector(0,-1){8}}
\put(64,12){\framebox(38,8){\shortstack[c]{Period turnover and\\political-culture state}}}
\put(64,16){\vector(-1,0){40}}
\put(83,20){\vector(0,1){8}}
\put(13,38){\makebox(0,0){\shortstack[c]{appointment, size, tenure,\\thresholds, councils, overrides}}}
\put(115,19){\makebox(0,0){\shortstack[c]{direct outputs, composites,\\uncertainty, provenance}}}
\put(64,5){\makebox(0,0){\shortstack[c]{Case generation, routing, decisions, compliance, reaction, and scoring are separate audit points.}}}
\end{picture}
\endgroup

\caption{Simulation pipeline and legislative-output bridge.}
\label{fig:model-pipeline}
\end{figure}

Each scenario creates justices or review actors from appointment method, court size, term and removal structure, independence level, accountability pressure, transparency level, coalition norm, rights priority, stability preference, review archetype, docket-control mechanism, review timing, cost profile, scenario kind, and review mechanism. The review process then models intake selection, emergency handling, recusals, panel or en banc review, voting thresholds, cross-checking agreement, constitutional-council screening, weak-form or suspended declarations, legislative response, explicit override, ombudsman access, public-defender participation, rights-impact statements, and opinion fragmentation. The reported campaigns use fixed weights rather than fitted parameters.

The weights have different epistemic status. Some inputs are institutional mappings, such as court size, route type, and whether a design uses weak-form declarations or merits follow-up. Some are source-observed benchmark ranges used only for diagnostic checks. Some are literature-informed directional priors: for example, higher political conflict should reduce compliance, and reason-giving should improve process transparency. Some are design-prior weights used to make synthetic mechanisms comparable under shared worlds. Finally, the directional score is a display convention, not a parameter estimate. The analysis below treats source-observed checks, synthetic mechanism movement, and normative score construction as separate claims.

The main model equations are intentionally simple enough to audit; the supplementary appendix states the fuller notation, composite-score equations, and route-mix construction. Review propensity is a weighted score:
\[
\begin{aligned}
R={}&A_s+.26T_r+.22S_c+.16L_c+.22C_p+.12I_c+.18E\\
&+.16O+.13D+.10P+.12B+.09L_a+.12P_i+.05G_r-.05U,
\end{aligned}
\]
where $A_s$ is structural access, $T_r$ rights threat, $S_c$ constitutional salience, $L_c$ lower-court conflict, $C_p$ certiorari or intake pressure, $I_c$ intercourt conflict, $E$ emergency status, $O$ ombudsman trigger, $D$ public-defender participation, $P$ pre-enactment review, $B$ abstract review, $L_a$ litigant capacity, $P_i$ public-interest support, $G_r$ government repeat-player advantage, and $U$ legal ambiguity. The case is reviewed when $R$ multiplied by jurisdictional access plus random docket noise exceeds the review threshold. Government repeat-player advantage can raise case salience and review probability while still lowering case-selection access and democratic-constitutionalism scores when it reflects asymmetric litigating capacity.

The review vote uses justice-level rights sensitivity, partisan pressure, legal ambiguity, lower-court signal, mandate restraint, and mechanism adjustments. Mechanism adjustments are deliberately directional rather than fitted: rights-impact statements slightly reduce unnecessary invalidation while increasing attention to high-rights-threat cases; ombudsman and public-defender participation increase access and rights salience; pre-enactment and abstract review increase salience screening; weak-form review lowers the threshold for a declaration without automatically invalidating the law.

Emergency orders are modeled separately from merits invalidation. The simulator tracks interim relief, reason-giving, vote disclosure, public disagreement, government applicants, government wins, and merits follow-up. This distinction matters because an institution can be restrained on the merits while still aggressive on emergency relief, or transparent in merits decisions while opaque on emergency orders.

Figure~\ref{fig:emergency-docket} illustrates why that separation matters. The current-style scenario generates high emergency relief and high emergency-process irregularity, while reason-giving and cross-checking designs can keep emergency relief available while sharply changing transparency and follow-up behavior.

\begin{figure}[hbt!]
\centering
% Auto-generated by paper/scripts/generate_figures.py
\begingroup
\setlength{\unitlength}{1mm}
\begin{picture}(128,72)
\scriptsize
\put(36.0,10){\color{black!15}\line(0,1){50}}
\put(36.0,6){\makebox(0,0){0}}
\put(55.5,10){\color{black!15}\line(0,1){50}}
\put(55.5,6){\makebox(0,0){0.25}}
\put(75.0,10){\color{black!15}\line(0,1){50}}
\put(75.0,6){\makebox(0,0){0.5}}
\put(94.5,10){\color{black!15}\line(0,1){50}}
\put(94.5,6){\makebox(0,0){0.75}}
\put(114.0,10){\color{black!15}\line(0,1){50}}
\put(114.0,6){\makebox(0,0){1}}
\put(36.0,10){\line(1,0){78.0}}
\put(36.0,10){\line(0,1){50}}
\put(36,64){\color{black}\rule{6mm}{1.2mm}}
\put(44,64){\makebox(0,0)[l]{shadow-docket abuse}}
\put(76,64){\color{black!55}\rule{6mm}{1.2mm}}
\put(84,64){\makebox(0,0)[l]{reason-giving}}
\put(36,60.5){\color{black!25}\rule{6mm}{1.2mm}}
\put(44,60.5){\makebox(0,0)[l]{merits follow-up}}
\put(32,56.0){\makebox(0,0)[r]{\color{red}CUR}}
\put(36.0,57.8){\color{black}\rule{28.5mm}{1.2mm}}
\put(36.0,56.0){\color{black!55}\rule{2.0mm}{1.2mm}}
\put(36.0,54.2){\color{black!25}\rule{0.0mm}{1.2mm}}
\put(32,48.0){\makebox(0,0)[r]{\color{black}REC}}
\put(36.0,49.8){\color{black}\rule{4.8mm}{1.2mm}}
\put(36.0,48.0){\color{black!55}\rule{54.4mm}{1.2mm}}
\put(36.0,46.2){\color{black!25}\rule{22.5mm}{1.2mm}}
\put(32,40.0){\makebox(0,0)[r]{\color{black}EMR}}
\put(36.0,41.8){\color{black}\rule{0.2mm}{1.2mm}}
\put(36.0,40.0){\color{black!55}\rule{67.7mm}{1.2mm}}
\put(36.0,38.2){\color{black!25}\rule{78.0mm}{1.2mm}}
\put(32,32.0){\makebox(0,0)[r]{\color{black}DUAL}}
\put(36.0,33.8){\color{black}\rule{0.2mm}{1.2mm}}
\put(36.0,32.0){\color{black!55}\rule{69.9mm}{1.2mm}}
\put(36.0,30.2){\color{black!25}\rule{78.0mm}{1.2mm}}
\put(32,24.0){\makebox(0,0)[r]{\color{black}COUNC}}
\put(36.0,25.8){\color{black}\rule{4.6mm}{1.2mm}}
\put(36.0,24.0){\color{black!55}\rule{58.3mm}{1.2mm}}
\put(36.0,22.2){\color{black!25}\rule{22.9mm}{1.2mm}}
\put(32,16.0){\makebox(0,0)[r]{\color{black}HYB}}
\put(36.0,17.8){\color{black}\rule{0.2mm}{1.2mm}}
\put(36.0,16.0){\color{black!55}\rule{68.2mm}{1.2mm}}
\put(36.0,14.2){\color{black!25}\rule{78.0mm}{1.2mm}}
\put(75,2){\makebox(0,0){rate}}
\end{picture}
\endgroup

\caption{Emergency-docket diagnostics in the baseline campaign. Higher reason-giving and merits follow-up are generally desirable; lower emergency-process irregularity is generally desirable.}
\label{fig:emergency-docket}
\end{figure}

A simple toy case clarifies the state transition. Suppose the docket contains a high-rights-burden election dispute with high urgency, weak vehicle quality, substantial lower-court conflict, and strong executive pressure. Under a fast emergency docket, the case can receive interim relief before merits review, producing rights-claimant success in the short run but also raising emergency-process irregularity when reasons, vote disclosure, or merits follow-up are absent. Under a written-reasons design, the same emergency relief is more transparent, so the direct opacity component falls, but the model still asks whether later merits review, compliance, and conflict improve. Under a no-relief-without-merits connection, the case is forced toward expedited merits handling; the expected emergency-irregularity score falls more directly, but delay, administrative load, and the probability of a denied urgent claimant can rise. This toy example shows why the paper treats the emergency findings as most interesting when they change downstream compliance, conflict, cost, or rights-claimant success, not merely when they lower a component that the rule itself defines.

\begin{table}[hbt!]
\centering
\caption{Mechanical implications versus diagnostic findings.}
\label{tab:mechanical-diagnostic}
\scriptsize
\setlength{\tabcolsep}{3pt}
\begin{tabular}{@{}p{0.24\linewidth}p{0.24\linewidth}p{0.42\linewidth}@{}}
\toprule
Modeled pattern & Mechanically implied? & What would make it diagnostically interesting \\
\midrule
Written emergency reasons reduce the opacity portion of emergency irregularity & Partly yes & Whether reasons also change merits follow-up, compliance, conflict, public-legitimacy proxy, or administrative cost. \\
No-relief-without-merits rules reduce merits-displacing emergency action & Partly yes & Whether urgent rights-claimant success is preserved and whether the downstream emergency effect falls without unacceptable delay. \\
Appointment or size reforms move emergency diagnostics less directly & No direct coding implication & Whether composition changes alter emergency outcomes after shared dockets, legal ambiguity, replacement timing, and public trust are held constant. \\
Weak-form or response-cycle mechanisms improve democratic constitutionalism & Not by definition & Whether legislative responses are timely and credible rather than merely relocating veto power outside courts. \\
\bottomrule
\end{tabular}
\end{table}

Runs are partitioned into review periods so term limits and removal or retention rules can change court composition over time. Decision outcomes also update an institutional reaction state: public trust, legislature-court conflict, court-curbing pressure, override pressure, amendment pressure, and compliance norms feed later cases in the same run. Enforcement is split into executive implementation, agency nonacquiescence, legislative reenactment, and local-government compliance. World assumptions include party fragmentation, governing coalition control, electoral time pressure, civil-society capacity, implementation capacity, and legal-tradition compatibility. For the validation campaign, real-world presets use country-year context modifiers; for synthetic campaigns, the same variables remain stress-test levers.

The main outcome metrics are listed in Table~\ref{tab:metrics}. The table separates direct procedural outputs, direct outcome outputs, constructed indices, and display-only reading aids. That distinction matters. Emergency reason-giving, vote disclosure, merits follow-up, weak-form declarations, legislative responses, and public-defender participation are direct model events. Rights-claimant success and invalidation are direct outcomes, but they are not identical to morally correct rights protection; a successful claimant can be strategically selected, anti-regulatory, or in tension with other rights holders. Legitimacy is therefore labeled as a modeled public-legitimacy proxy rather than an opinion-survey measure, and the legacy \texttt{shadowDocketAbuse} report column is interpreted as emergency-process irregularity rather than a legal conclusion that a specific order is abusive. Constitutional conflict is likewise a composite index covering legal, political, federalism, and implementation conflict.

The new constitutional-design diagnostics are synthetic constructs. Veto-relocation risk rises when a design shifts blocking power into courts, councils, cross-checks, or upstream admissibility screens without a rights-threat justification or legislative response. Legal-transplant feasibility combines scenario transplant score, political-culture fit, trust, inverse institutional demand, legal-tradition compatibility, implementation capacity, civil-society capacity, and inverse political mismatch. Political-culture sensitivity increases when a mechanism depends heavily on cooperative legislative response, high public trust, low polarization, low fragmentation, or strong compliance norms. Democratic constitutionalism combines the values most central to the paper's question:
\[
\begin{aligned}
D_c={}&.11S+.18R_p+.12L+.15D_r+.09C+.08(1-K)+.07(1-V)+.05F\\
&+.03(1-P_c)+.03Q+.03A_c+.03(1-G_r)+.03I_m,
\end{aligned}
\]
where $S$ is legal stability, $R_p$ rights protection, $L$ the modeled public-legitimacy proxy, $D_r$ democratic responsiveness, $C$ compliance, $K$ the constitutional-conflict index, $V$ veto-relocation risk, $F$ transplant feasibility, $P_c$ political-culture sensitivity, $Q$ legislative-response credibility, $A_c$ case-selection access, $G_r$ government repeat-player advantage, and $I_m$ implementation capacity. The directional score remains only a display aid. It combines directional metrics after inverting lower-is-better values; it should not be treated as a final welfare function.

\begin{table}[hbt!]
\centering
\caption{Core metric families.}
\label{tab:metrics}
\begin{tabular}{@{}p{0.23\linewidth}p{0.50\linewidth}p{0.13\linewidth}@{}}
\toprule
Metric family & Examples & Direction \\
\midrule
Legal performance & legal stability, rights protection, reversal rate, merits invalidation rate & mixed \\
Political alignment & partisan alignment, democratic responsiveness, independence/accountability balance & mixed \\
Democratic constitutionalism & democratic constitutionalism, legislative-response credibility, case-selection access, government repeat-player advantage, implementation capacity, veto-relocation risk, transplant feasibility, political-culture sensitivity & mixed \\
Noncourt mechanisms & weak-form declarations, suspended declarations, legislative response cycles, rights-impact statements, ombudsman triggers, public-defender participation, pre-enactment review, abstract review & diagnostic \\
Supranational route mix & preliminary references, appeal routes, direct actions & diagnostic \\
Emergency docket & emergency relief, reason-giving, vote disclosure, public disagreement, merits follow-up & mixed \\
Legitimacy and conflict & modeled public-legitimacy proxy, public trust input, constitutional-conflict index, court-curbing pressure, amendment pressure & mixed \\
Pipeline and intake & review rate, intake acceptance, case-selection access, litigant capacity, public-interest support, lower-court depth, state-federal tension, intercourt conflict & diagnostic \\
Coalitions and recusals & concurrence fragmentation, dissent intensity, recusal rate, en banc rate, cross-check rate & mixed \\
Compliance & compliance, defiance, workarounds, executive implementation, agency nonacquiescence, reenactment, local compliance & mixed \\
Institutional cost & administrative load, direct court cost, upstream screening cost, capacity strain, delay, implementation complexity & lower \\
\bottomrule
\end{tabular}
\end{table}

\subsection{Directional score construction}

The directional score is an unweighted diagnostic summary, not an estimated social welfare function. Let $H$ be the set of higher-is-better metrics and $L$ the set of lower-is-better metrics included in the display score. The score is
\[
S=\frac{1}{35}\left(\sum_{m\in H}m+\sum_{m\in L}(1-m)\right).
\]
The higher-is-better set includes legal stability, rights protection, modeled public-legitimacy proxy, democratic responsiveness, legislative-response credibility, timely legislative response, case-selection access, implementation capacity, democratic constitutionalism, independence/accountability balance, compliance, executive implementation, local-government compliance, emergency reason-giving, emergency vote disclosure, legal-transplant feasibility, legislative response, and rights-impact statements. The lower-is-better set includes partisan alignment, emergency-process irregularity, reversal rate, emergency relief rate, merits invalidation rate, constitutional-conflict index, government repeat-player advantage, veto-relocation risk, political-culture sensitivity, defiance, workarounds, agency nonacquiescence, legislative reenactment, emergency public disagreement, legislative-response delay, administrative load, and total institutional cost. Diagnostic metrics such as review rate, recusal rate, override rate, weak-form declaration rate, abstract-review rate, lower-court depth, replacement pressure, and public trust remain visible in reports but do not enter $S$.

Small differences in $S$ are not substantively meaningful. The tables and figures therefore use the score only as a compact index for orienting the reader to clusters of designs. Designs with overlapping bootstrap intervals or score differences below about one percentage point should be treated as tied unless an individual mechanism diagnostic explains a stable difference. The main results privilege multi-objective comparisons--emergency irregularity, rights-claimant success, democratic constitutionalism, public-legitimacy proxy, veto-relocation risk, transplant feasibility, political-culture sensitivity, and cost--over ordinal rankings by this score.

\section{Calibration, Uncertainty, and Replication}
\label{sec:calibration}

The model is not empirically validated in the sense of fitting a historical court. It uses documented external ranges in two different ways. Source-specific diagnostic checks require a citable source, nonzero denominator, direct target analogue, and no synthesis or normalized-cost construction. Contextual benchmarks, public-trust ranges, doctrine-mix proxies without stored denominators, and institutional-cost ranges are used only as stress-test assumptions. They help discipline the comparative setup, but they are not calibration evidence and are not counted in validation-style summaries.

Calibration targets identify the court, time period, target measure, target range, observed value when available, sample size, unit, method, reliability, source, validation-use flag, and construction note. The current targets include U.S. Supreme Court merits-docket doctrine shares from the Supreme Court Database \citep{scdb2025,scdbIssueCodebook}, emergency-docket context \citep{scotusblogInterim2024}, comparative court intake and throughput proxies, public-trust context \citep{gallupCourtApproval2024}, official-source intake or response rows for Canada, France, the United Kingdom, the European Court of Human Rights, and the CJEU, plus normalized institutional cost profiles. The United Kingdom permission-to-appeal row is also mirrored as a case-selection-access proxy; that row is used only to discipline the access diagnostic and should not be read as a direct measure of legal aid, counsel quality, filing cost, standing, or public-interest litigation capacity. Under the source-specific rule, validation-counted rows remain narrow: they require a documented source, denominator, and direct simulator analogue. The comparative presets for Germany, South Africa, and the U.S. merits docket are therefore still presented as structured real-world presets and stress tests, not as empirical validations.

The calibration workflow also records promising but not yet validation-ready findings. Some observations have been incorporated after source verification; others remain excluded because they are secondary-source findings, lack a direct simulator analogue, or need a coding-rule audit. These excluded observations include leads on Canadian Charter legislative sequels, conditional French QPC remedy rates, comparative pre-enactment review frequencies, and transplant-feasibility indicators from public-law and governance datasets. This separation is meant to prevent a common error in simulation papers: allowing source-discovery notes to masquerade as calibration.

% Auto-generated by paper/scripts/generate_tables.py
\begin{table}[hbt!]
\centering
\caption{Examples of external calibration targets.}
\label{tab:calibration-targets}
\scriptsize
\setlength{\tabcolsep}{3pt}
\begin{tabular}{@{}p{0.26\linewidth}p{0.27\linewidth}rrrl@{}}
\toprule
Court and period & Target & Observed & Range & $n$ & Reliability \\
\midrule
U.S. Supreme Court (2000-2024 terms) & Speech docket share & 0.060 & 0.039--0.080 & -- & high \\
U.S. Supreme Court (2000-2024 terms) & Civil-rights and privacy docket share & 0.164 & 0.137--0.191 & -- & high \\
U.S. Supreme Court (2024-2025 emergency docket) & Substantive emergency application relief rate & 0.440 & 0.310--0.460 & 43 & medium \\
U.S. Supreme Court (2024-2025 emergency docket) & Written explanation share & 0.279 & 0.200--0.360 & 43 & medium \\
U.S. Supreme Court (2024 public-opinion year) & Public court trust and approval & 0.420 & 0.350--0.490 & -- & medium \\
German Federal Constitutional Court (2024) & Constitutional complaint success and admission proxy & 0.009 & 0.006--0.012 & 4640 & medium \\
Supreme Court of Canada (recent annual average) & Leave application grant rate & 0.089 & 0.075--0.105 & 524 & high \\
French Constitutional Council (QPC recent period) & QPC invalidation rate & 0.315 & 0.260--0.360 & 67 & medium \\
Constitutional Court of South Africa (recent annual average) & Petition-to-judgment throughput proxy & 0.141 & 0.110--0.170 & 355 & medium \\
\bottomrule
\end{tabular}
\begin{minipage}{0.96\linewidth}
\footnotesize Notes: The table reports selected calibration targets from the repository's CSV files. Ranges are target intervals, not estimated model effects. Source names, URLs, notes, and validation flags are retained in the replication files.
\end{minipage}
\end{table}


The reported campaigns include case-level outcome records and bootstrap uncertainty bands for aggregate and segment diagnostics. The bootstrap resamples entire generated-world run blocks rather than isolated cases, which preserves within-run dependence while still using the raw outcome records. Composition and calibration intervals remain conservative denominator-based checks where no case-level analogue exists.

\subsection{Validation-style diagnostics}

The validation campaign runs real-world scenario presets against a shared benchmark docket and compares the resulting diagnostics to documented calibration ranges where a direct analogue exists. Table~\ref{tab:validation-summary} reports the count of mapped targets that fall within range and the largest remaining miss for each profile. These checks are intentionally diagnostic. A miss may indicate that the simulator abstracts away a case-flow mechanism, that a target is contextual rather than directly comparable, or that the preset should be refined; it is not treated as a failed prediction.

% Auto-generated by paper/scripts/generate_tables.py
\begin{table}[hbt!]
\centering
\caption{Validation-style checks against documented target ranges.}
\label{tab:validation-summary}
\scriptsize
\setlength{\tabcolsep}{4pt}
\begin{tabular}{@{}p{0.22\linewidth}p{0.25\linewidth}rrp{0.27\linewidth}@{}}
\toprule
Profile & Scenario preset & In range & Median gap & Largest miss \\
\midrule
U.S. Supreme Court & U.S. Supreme Court benchmark & 1/7 & 0.066 & Emergency-powers merits share \\
U.S. emergency docket & U.S. Supreme Court benchmark & 0/4 & 0.176 & Substantive emergency application relief rate \\
Germany BVerfG & German Federal Constitutional Court & 1/1 & 0.000 & none \\
Canada SCC & Canadian Supreme Court with override context & 1/1 & 0.000 & none \\
France Conseil & French Constitutional Council & 1/1 & 0.000 & none \\
South Africa ConstCourt & South African Constitutional Court & 0/1 & 0.830 & Petition-to-judgment throughput proxy \\
\bottomrule
\end{tabular}
\begin{minipage}{0.96\linewidth}
\footnotesize Notes: The validation campaign runs real-world scenario presets against a shared benchmark docket. The checks count only source-backed target rows marked for validation in the calibration source matrix; synthesis, public-trust, and normalized-cost rows are excluded. They are external diagnostics, not fitted validation estimates.
\end{minipage}
\end{table}


Table~\ref{tab:validation-miss-interpretation} makes the largest misses explicit. Intake-rate misses usually point to a denominator problem: the simulator's generated filing pool and screeners do not yet reproduce a court's annual leave, petition, or admissibility universe. Case-selection-access misses are even narrower because the current promoted row is a permission-to-appeal grant-rate proxy rather than a full access measure. Emergency-relief misses point to procedure-specific assumptions about applicant mix, urgency, interim relief, public reasons, and merits follow-through. Weak-form response misses point to the quality and timing of legislative response rather than to the court's merits behavior. These categories prevent the external checks from being read as a single pass/fail validation score.

% Auto-generated by paper/scripts/generate_tables.py
\begin{table}[hbt!]
\centering
\caption{Calibration roadmap from the largest source-range misses.}
\label{tab:validation-miss-interpretation}
\scriptsize
\setlength{\tabcolsep}{3pt}
\begin{tabular}{@{}L{0.24\linewidth}L{0.18\linewidth}L{0.21\linewidth}L{0.29\linewidth}@{}}
\toprule
Target & Model vs. source & Category & Interpretation and next step \\
\midrule
Court of Justice of the European Union: Full compliance with rule-of-law rulings & 0.642 vs. 0.534--0.632; gap 0.010 & source-range check & The preset misses this source-range check. Next: Inspect whether the source denominator, simulator metric, and scenario preset are directly comparable. \\
French Constitutional Council: QPC nonconformity rate & 0.299 vs. 0.305--0.324; gap 0.006 & merits-outcome mechanism & The merits-outcome rate is close but outside the source interval for this profile. Next: Retune doctrine severity, referral screening, and invalidation-threshold parameters only after preserving cross-profile comparability. \\
\bottomrule
\end{tabular}
\begin{minipage}{0.96\linewidth}
\footnotesize Notes: The table reports source-range misses when validation-counted checks fall outside documented ranges; when none do, it records the no-current-miss state. The full generated source-range report is retained in \texttt{reports/constitutional-review-validation-v1-misses.csv} and \texttt{.md}.
\end{minipage}
\end{table}


\subsection{What the diagnostics still miss}

The validation diagnostics are intentionally better at finding missing mechanisms than at confirming model fit. The most important current misses are pipeline misses: generated filing pools, leave filters, lower-court conflict formation, repeat-player selection, and public-interest support are still stylized. Compliance and enforcement diagnostics are downstream model outputs, not source-fitted country estimates. Emergency-docket diagnostics distinguish reasons, vote disclosure, public disagreement, government applicants, merits follow-up, and relief, but they do not yet use a docket-level dataset of applications, responses, referral decisions, and subsequent merits disposition. These gaps are why the paper separates empirical claims, synthetic findings, and speculative design recommendations. The empirical claims are limited to the existence and construction of source-bearing target rows. The synthetic findings describe how fixed model mechanisms move under shared simulated dockets. The design recommendations are conditional interpretations of those synthetic movements, not claims that a named court would behave that way.

The narrowness of Tables~\ref{tab:validation-summary} and \ref{tab:validation-miss-interpretation} is itself a result. Comparative rows are excluded from the validation count when they rely on synthesis labels, missing denominators, public-trust proxies, unsupported target keys, or normalized-cost assumptions. Source-counted rows for Canada, France, the United Kingdom, the ECHR, the CJEU, and the U.S. emergency docket are still narrow, source-specific diagnostics. That leaves a smaller validation-style surface and makes the limitation explicit: the framework is ready for comparative design stress testing, not for estimating the behavior of named courts. Empirical extensions should replace each contextual row with a source-specific dataset before using the word validation for country profiles.

Replication materials separate source code, model inputs, generated reports, manuscript source, and provenance records. The public replication deposit should include source code, model inputs, aggregate reports, full case-level and interval outputs or scripts that regenerate them, and instructions for reproducing the campaign outputs.

\section{Comparative Experiments}
\label{sec:experiments}

The main campaign compares court-variant scenarios and synthetic mechanism scenarios across five stress cases: baseline, polarized appointments, rights-threat surge, emergency-docket stress, and low-trust conflict. A paired-input campaign then compares the same designs against a synthetic baseline and four legislative-output docket modes: all imported rows, high-capture rows, high-volatility rows, and low-mandate rows. The sensitivity campaign reruns the scenario catalog under high and low assumptions for emergency pressure, appointment polarization, rights-threat rate, public trust, and legislature-court conflict.

Table~\ref{tab:baseline-results} reports the baseline estimates that anchor the graphical summaries. The table puts direct and constructed diagnostics before the display score so that the reader can evaluate emergency-process irregularity, rights-claimant success, modeled legitimacy, veto relocation, feasibility, and cost without treating the score as a ranking device.

% Auto-generated by paper/scripts/generate_tables.py
\begin{table}[hbt!]
\centering
\caption{Court-architecture stress-test results.}
\label{tab:baseline-results}
\scriptsize
\setlength{\tabcolsep}{2pt}
\begin{tabular}{@{}lL{0.26\linewidth}rrrrrrr@{}}
\toprule
Code & Architecture & Dem. const. & Rights & Emerg. irr. $\downarrow$ & Legit. proxy & Veto $\downarrow$ & Cost $\downarrow$ & Score aid \\
\midrule
CUR & Current-style discretionary apex court & 0.571 & 0.687 & 0.366 & 0.582 & 0.529 & 0.450 & 0.525 \\
18Y & Eighteen-year terms & 0.578 & 0.695 & 0.359 & 0.592 & 0.520 & 0.455 & 0.535 \\
15J & Fifteen-justice commission & 0.643 & 0.739 & 0.062 & 0.708 & 0.436 & 0.470 & 0.615 \\
SUP & Supermajority invalidation rule & 0.638 & 0.736 & 0.064 & 0.684 & 0.446 & 0.457 & 0.604 \\
REC & Strict recusal court & 0.648 & 0.739 & 0.063 & 0.726 & 0.434 & 0.476 & 0.622 \\
EMR & Reasoned emergency review & 0.649 & 0.740 & 0.003 & 0.740 & 0.446 & 0.460 & 0.625 \\
PAN & Panel and en banc routing & 0.643 & 0.739 & 0.062 & 0.720 & 0.433 & 0.511 & 0.609 \\
DUAL & Dual cross-checking courts & 0.650 & 0.747 & 0.003 & 0.751 & 0.596 & 0.589 & 0.628 \\
COUNC & Constitutional council & 0.652 & 0.741 & 0.058 & 0.753 & 0.662 & 0.533 & 0.632 \\
OVR & Legislative override court & 0.637 & 0.733 & 0.064 & 0.700 & 0.435 & 0.466 & 0.592 \\
RET & Retention-accountability court & 0.656 & 0.729 & 0.057 & 0.723 & 0.407 & 0.482 & 0.622 \\
HYB & Independence-accountability hybrid & 0.663 & 0.746 & 0.002 & 0.754 & 0.433 & 0.542 & 0.640 \\
\bottomrule
\end{tabular}
\begin{minipage}{0.96\linewidth}
\footnotesize Notes: Values are baseline campaign estimates for court-architecture scenarios only; mechanism scenarios are separated in Table~\ref{tab:mechanism-results}. Lower emergency irregularity indicates less opaque or merits-displacing emergency intervention. The legitimacy column is a modeled public-legitimacy proxy. The score aid is the equally weighted directional display score; its 95\% run-block bootstrap interval is reported in the replication files.
\end{minipage}
\end{table}


Table~\ref{tab:mechanism-results} isolates the noncourt mechanisms. These estimates are synthetic findings only: they compare mechanisms under common simulated dockets and should not be read as evidence that any country using a similar device would obtain the displayed result.

% Auto-generated by paper/scripts/generate_tables.py
\begin{table}[hbt!]
\centering
\caption{Mechanism activation and downstream propagation effects.}
\label{tab:mechanism-results}
\scriptsize
\setlength{\tabcolsep}{2pt}
\begin{tabular}{@{}lL{0.18\linewidth}L{0.22\linewidth}rrrrrrr@{}}
\toprule
Code & Mechanism & Direct activation & Rights & Compliance & Conflict $\downarrow$ & Legit. & Cost $\downarrow$ & Veto $\downarrow$ & Dem. const. \\
\midrule
WFR & Weak-form review & Weak 0.357; Resp. 0.376 & 0.718 & 0.802 & 0.208 & 0.754 & 0.359 & 0.177 & 0.709 \\
SUS & Suspended declaration & Susp. 0.442; Resp. 0.451 & 0.724 & 0.596 & 0.246 & 0.761 & 0.383 & 0.207 & 0.680 \\
OVC & Override clause & Override 0.000; Resp. 0.261 & 0.738 & 0.605 & 0.221 & 0.737 & 0.425 & 0.355 & 0.655 \\
PRE & Pre-enactment review & RIS 0.381; Pre 0.833; Abs 0.872 & 0.766 & 0.683 & 0.234 & 0.755 & 0.437 & 0.511 & 0.667 \\
ABS & Abstract review tribunal & RIS 0.381; Pre 0.833; Abs 0.872 & 0.770 & 0.603 & 0.234 & 0.747 & 0.436 & 0.312 & 0.665 \\
OMB & Ombudsman-triggered review & Omb. 0.263 & 0.736 & 0.533 & 0.252 & 0.729 & 0.516 & 0.525 & 0.632 \\
DEF & Constitutional public defender & Def. 0.306 & 0.773 & 0.537 & 0.240 & 0.741 & 0.518 & 0.514 & 0.646 \\
RIS & Rights-impact statement & RIS 1.000; Pre 0.833; Abs 0.872 & 0.805 & 0.720 & 0.231 & 0.773 & 0.437 & 0.502 & 0.691 \\
MLR & Mandatory legislative response & Resp. 0.390 & 0.717 & 0.665 & 0.197 & 0.782 & 0.394 & 0.224 & 0.680 \\
\bottomrule
\end{tabular}
\begin{minipage}{0.96\linewidth}
\footnotesize Notes: Direct activation reports implementation-facing outputs; the remaining columns report downstream propagation effects. Front-end rows are bundle/routing scenarios rather than isolated toggles, so RIS, Pre, and Abs may all appear when a design routes review through multiple upstream screens. Lower conflict, cost, and veto-relocation risk are better; higher rights, compliance, legitimacy, and democratic constitutionalism are better.
\end{minipage}
\end{table}


Figure~\ref{fig:paired-import} shows the paired-input campaign as a design stress test. Legislative-output profiles can change the relative diagnostic score of each design, which is why cross-institutional stress testing is useful.

\begin{figure}[hbt!]
\centering
% Auto-generated by paper/scripts/generate_figures.py
\begingroup
\setlength{\unitlength}{1mm}
\begin{picture}(128,88)
\scriptsize
\put(24.0,14.0){\color{black!15}\line(1,0){82.0}}
\put(20.0,14.0){\makebox(0,0)[r]{0.50}}
\put(24.0,31.1){\color{black!15}\line(1,0){82.0}}
\put(20.0,31.1){\makebox(0,0)[r]{0.60}}
\put(24.0,48.1){\color{black!15}\line(1,0){82.0}}
\put(20.0,48.1){\makebox(0,0)[r]{0.70}}
\put(24.0,65.2){\color{black!15}\line(1,0){82.0}}
\put(20.0,65.2){\makebox(0,0)[r]{0.80}}
\put(24.0,14.0){\color{black!12}\line(0,1){58.0}}
\put(24.0,10.0){\makebox(0,0){SB}}
\put(44.5,14.0){\color{black!12}\line(0,1){58.0}}
\put(44.5,10.0){\makebox(0,0){ALL}}
\put(65.0,14.0){\color{black!12}\line(0,1){58.0}}
\put(65.0,10.0){\makebox(0,0){CAP}}
\put(85.5,14.0){\color{black!12}\line(0,1){58.0}}
\put(85.5,10.0){\makebox(0,0){VOL}}
\put(106.0,14.0){\color{black!12}\line(0,1){58.0}}
\put(106.0,10.0){\makebox(0,0){MAND}}
\put(24.0,14.0){\line(1,0){82.0}}
\put(24.0,14.0){\line(0,1){58.0}}
\linethickness{0.35mm}
{\color{black}\qbezier[22](24.0,13.3)(34.2,24.7)(44.5,36.2)}
{\color{black}\qbezier[22](44.5,36.2)(54.8,33.1)(65.0,30.0)}
{\color{black}\qbezier[22](65.0,30.0)(75.2,26.6)(85.5,23.2)}
{\color{black}\qbezier[22](85.5,23.2)(95.8,24.4)(106.0,25.6)}
\linethickness{0.25mm}
\put(24.0,13.3){\makebox(0,0){\color{black}\rule{1.5mm}{1.5mm}}}
\put(44.5,36.2){\makebox(0,0){\color{black}\rule{1.5mm}{1.5mm}}}
\put(65.0,30.0){\makebox(0,0){\color{black}\rule{1.5mm}{1.5mm}}}
\put(85.5,23.2){\makebox(0,0){\color{black}\rule{1.5mm}{1.5mm}}}
\put(106.0,25.6){\makebox(0,0){\color{black}\rule{1.5mm}{1.5mm}}}
\linethickness{0.35mm}
{\color{black!70}\qbezier[22](24.0,32.3)(34.2,40.7)(44.5,49.1)}
{\color{black!70}\qbezier[22](44.5,49.1)(54.8,48.9)(65.0,48.6)}
{\color{black!70}\qbezier[22](65.0,48.6)(75.2,43.7)(85.5,38.7)}
{\color{black!70}\qbezier[22](85.5,38.7)(95.8,39.7)(106.0,40.6)}
\linethickness{0.25mm}
\put(24.0,32.3){\makebox(0,0){\color{black!70}\rule{1.5mm}{1.5mm}}}
\put(44.5,49.1){\makebox(0,0){\color{black!70}\rule{1.5mm}{1.5mm}}}
\put(65.0,48.6){\makebox(0,0){\color{black!70}\rule{1.5mm}{1.5mm}}}
\put(85.5,38.7){\makebox(0,0){\color{black!70}\rule{1.5mm}{1.5mm}}}
\put(106.0,40.6){\makebox(0,0){\color{black!70}\rule{1.5mm}{1.5mm}}}
\linethickness{0.35mm}
{\color{black!45}\qbezier[22](24.0,31.6)(34.2,39.3)(44.5,47.1)}
{\color{black!45}\qbezier[22](44.5,47.1)(54.8,42.3)(65.0,37.5)}
{\color{black!45}\qbezier[22](65.0,37.5)(75.2,38.2)(85.5,38.9)}
{\color{black!45}\qbezier[22](85.5,38.9)(95.8,36.3)(106.0,33.8)}
\linethickness{0.25mm}
\put(24.0,31.6){\makebox(0,0){\color{black!45}\rule{1.5mm}{1.5mm}}}
\put(44.5,47.1){\makebox(0,0){\color{black!45}\rule{1.5mm}{1.5mm}}}
\put(65.0,37.5){\makebox(0,0){\color{black!45}\rule{1.5mm}{1.5mm}}}
\put(85.5,38.9){\makebox(0,0){\color{black!45}\rule{1.5mm}{1.5mm}}}
\put(106.0,33.8){\makebox(0,0){\color{black!45}\rule{1.5mm}{1.5mm}}}
\linethickness{0.35mm}
{\color{black!30}\qbezier[22](24.0,29.7)(34.2,38.4)(44.5,47.1)}
{\color{black!30}\qbezier[22](44.5,47.1)(54.8,46.9)(65.0,46.8)}
{\color{black!30}\qbezier[22](65.0,46.8)(75.2,43.9)(85.5,41.1)}
{\color{black!30}\qbezier[22](85.5,41.1)(95.8,42.4)(106.0,43.7)}
\linethickness{0.25mm}
\put(24.0,29.7){\makebox(0,0){\color{black!30}\rule{1.5mm}{1.5mm}}}
\put(44.5,47.1){\makebox(0,0){\color{black!30}\rule{1.5mm}{1.5mm}}}
\put(65.0,46.8){\makebox(0,0){\color{black!30}\rule{1.5mm}{1.5mm}}}
\put(85.5,41.1){\makebox(0,0){\color{black!30}\rule{1.5mm}{1.5mm}}}
\put(106.0,43.7){\makebox(0,0){\color{black!30}\rule{1.5mm}{1.5mm}}}
\linethickness{0.35mm}
{\color{black!55}\qbezier[22](24.0,33.8)(34.2,41.8)(44.5,49.8)}
{\color{black!55}\qbezier[22](44.5,49.8)(54.8,48.7)(65.0,47.6)}
{\color{black!55}\qbezier[22](65.0,47.6)(75.2,44.5)(85.5,41.3)}
{\color{black!55}\qbezier[22](85.5,41.3)(95.8,40.8)(106.0,40.3)}
\linethickness{0.25mm}
\put(24.0,33.8){\makebox(0,0){\color{black!55}\rule{1.5mm}{1.5mm}}}
\put(44.5,49.8){\makebox(0,0){\color{black!55}\rule{1.5mm}{1.5mm}}}
\put(65.0,47.6){\makebox(0,0){\color{black!55}\rule{1.5mm}{1.5mm}}}
\put(85.5,41.3){\makebox(0,0){\color{black!55}\rule{1.5mm}{1.5mm}}}
\put(106.0,40.3){\makebox(0,0){\color{black!55}\rule{1.5mm}{1.5mm}}}
\linethickness{0.35mm}
{\color{black!40}\qbezier[22](24.0,39.9)(34.2,45.8)(44.5,51.7)}
{\color{black!40}\qbezier[22](44.5,51.7)(54.8,50.1)(65.0,48.5)}
{\color{black!40}\qbezier[22](65.0,48.5)(75.2,48.3)(85.5,48.1)}
{\color{black!40}\qbezier[22](85.5,48.1)(95.8,49.3)(106.0,50.5)}
\linethickness{0.25mm}
\put(24.0,39.9){\makebox(0,0){\color{black!40}\rule{1.5mm}{1.5mm}}}
\put(44.5,51.7){\makebox(0,0){\color{black!40}\rule{1.5mm}{1.5mm}}}
\put(65.0,48.5){\makebox(0,0){\color{black!40}\rule{1.5mm}{1.5mm}}}
\put(85.5,48.1){\makebox(0,0){\color{black!40}\rule{1.5mm}{1.5mm}}}
\put(106.0,50.5){\makebox(0,0){\color{black!40}\rule{1.5mm}{1.5mm}}}
\linethickness{0.35mm}
{\color{black!60}\qbezier[22](24.0,36.0)(34.2,42.7)(44.5,49.5)}
{\color{black!60}\qbezier[22](44.5,49.5)(54.8,49.3)(65.0,49.1)}
{\color{black!60}\qbezier[22](65.0,49.1)(75.2,46.1)(85.5,43.0)}
{\color{black!60}\qbezier[22](85.5,43.0)(95.8,40.1)(106.0,37.2)}
\linethickness{0.25mm}
\put(24.0,36.0){\makebox(0,0){\color{black!60}\rule{1.5mm}{1.5mm}}}
\put(44.5,49.5){\makebox(0,0){\color{black!60}\rule{1.5mm}{1.5mm}}}
\put(65.0,49.1){\makebox(0,0){\color{black!60}\rule{1.5mm}{1.5mm}}}
\put(85.5,43.0){\makebox(0,0){\color{black!60}\rule{1.5mm}{1.5mm}}}
\put(106.0,37.2){\makebox(0,0){\color{black!60}\rule{1.5mm}{1.5mm}}}
{\color{black!55}\qbezier[10](107.1,25.6)(108.0,25.1)(109.0,24.6)}
\put(110.0,24.6){\makebox(0,0)[l]{\color{black}CUR*}}
{\color{black!55}\qbezier[10](107.1,40.6)(108.0,45.9)(109.0,51.2)}
\put(110.0,51.2){\makebox(0,0)[l]{\color{black}EMR}}
{\color{black!55}\qbezier[10](107.1,33.8)(108.0,37.0)(109.0,40.2)}
\put(110.0,40.2){\makebox(0,0)[l]{\color{black}DUAL}}
{\color{black!55}\qbezier[10](107.1,43.7)(108.0,50.2)(109.0,56.7)}
\put(110.0,56.7){\makebox(0,0)[l]{\color{black}COUNC}}
{\color{black!55}\qbezier[10](107.1,40.3)(108.0,43.0)(109.0,45.7)}
\put(110.0,45.7){\makebox(0,0)[l]{\color{black}HYB}}
{\color{black!55}\qbezier[10](107.1,50.5)(108.0,56.4)(109.0,62.2)}
\put(110.0,62.2){\makebox(0,0)[l]{\color{black}WFR}}
{\color{black!55}\qbezier[10](107.1,37.2)(108.0,36.0)(109.0,34.7)}
\put(110.0,34.7){\makebox(0,0)[l]{\color{black}MLR}}
\put(65,4){\makebox(0,0){docket source}}
\put(12,75){\makebox(0,0){score}}
\put(65,80){\makebox(0,0){SB generated; ALL all imports; CAP high capture; VOL high volatility; MAND low mandate}}
\end{picture}
\endgroup

\caption{Directional diagnostic scores under synthetic and legislative-output docket sources. The score is a compact display aid, not a welfare function. CUR* marks the current-style scenario.}
\label{fig:paired-import}
\end{figure}

Four patterns emerge from the experiments. First, emergency-process reforms produce the clearest movement on the paper's central diagnostic. Reasoned emergency panels and merits follow-up reduce emergency-process irregularity by moving more emergency matters into reasoned or merits-connected review. That direct movement is partly mechanical, so the relevant question is whether these reforms preserve rights-claimant success and improve downstream legitimacy, compliance, or conflict. Second, appointment, size, and selection reforms mostly act indirectly. They can change legitimacy and alignment diagnostics through replacement dynamics, but they do not discipline emergency procedure as directly as written-reasons or merits-follow-through rules. Third, cross-checking and council designs tend to improve stability and lower reversal rates, but they do so by increasing administrative load, upstream screening, and implementation complexity. Fourth, term limits and appointment systems matter most when linked to dynamic replacement over time; their effects are muted if composition is treated as fixed at initial court creation.

The mechanism results add a fourth pattern. Dialogic and front-end mechanisms can improve democratic-constitutionalism scores only when they also reduce veto-relocation risk or produce a modeled legislative response. Weak-form review without legislative uptake is less attractive in the simulation than weak-form or suspended-declaration designs that trigger response cycles. Rights-impact statements and pre-enactment review lower downstream conflict in some worlds, but they can also relocate veto power into front-end screening if transparency and political-culture fit are weak.

The negative results are equally important. The highest-scoring designs are not costless reforms, and the validation checks do not support a claim that any preset reproduces a real court. Close score differences should be treated as ties unless they are accompanied by clear and stable movement in a mechanism diagnostic. The current-style design remains administratively lean, which is a real institutional advantage, but its emergency-docket and modeled-legitimacy diagnostics are weaker in the baseline campaign. Council and cross-checking designs improve several legal-performance measures in the model, but they do so by relocating work to screening institutions or duplicate review structures. The useful output is therefore a tradeoff map, not a rank ordering.

These patterns should be read as comparative diagnostics rather than findings about any existing court. The model is useful because all scenarios face the same generated worlds and dockets. The relevant comparison is not whether a simulated current-style court matches a historical court. It is whether changing a design feature while holding the rest of the simulated environment constant exposes a tradeoff that constitutional designers should take seriously.

\section{Implications for Institutional Design}
\label{sec:implications}

The simulation framework points toward a bundle-based account of constitutional review. Independence rules do not have a single effect; they interact with docket control, voting thresholds, review timing, emergency procedure, and enforcement channels. Accountability mechanisms can improve democratic responsiveness but also increase partisan alignment or weaken rights protection if they operate through short renewable terms or retention pressure. Emergency reforms are especially sensitive to whether the reform only demands speed, only demands reasons, or requires a structured merits follow-up pathway.

The model also makes administrative externalities visible. A small council may look inexpensive if only the top body is counted, but referral-gated review shifts costs into lower courts, political actors, or admissibility filters. A dual-court design may increase legitimacy through cross-checking but raise delay and capacity strain. A declaration-oriented model may improve democratic responsiveness while reducing judicial finality. These are not purely technical details; they are constitutional design consequences.

Figure~\ref{fig:cost-score} displays the same point graphically. The current-style court is low cost but has a lower directional diagnostic score in the baseline campaign, while cross-checking, council, and hybrid designs move upward at the price of higher institutional cost. This is the kind of tradeoff that can disappear inside a single aggregate score.

\begin{inlinefigureblock}
% Auto-generated by paper/scripts/generate_figures.py
\begingroup
\setlength{\unitlength}{1mm}
\begin{picture}(128,86)
\scriptsize
\put(33.5,15.0){\color{black!15}\line(0,1){58.0}}
\put(33.5,11.0){\makebox(0,0){0.30}}
\put(56.5,15.0){\color{black!15}\line(0,1){58.0}}
\put(56.5,11.0){\makebox(0,0){0.40}}
\put(79.5,15.0){\color{black!15}\line(0,1){58.0}}
\put(79.5,11.0){\makebox(0,0){0.50}}
\put(102.5,15.0){\color{black!15}\line(0,1){58.0}}
\put(102.5,11.0){\makebox(0,0){0.60}}
\put(22.0,15.0){\color{black!15}\line(1,0){92.0}}
\put(18.0,15.0){\makebox(0,0)[r]{0.55}}
\put(22.0,29.5){\color{black!15}\line(1,0){92.0}}
\put(18.0,29.5){\makebox(0,0)[r]{0.60}}
\put(22.0,44.0){\color{black!15}\line(1,0){92.0}}
\put(18.0,44.0){\makebox(0,0)[r]{0.65}}
\put(22.0,58.5){\color{black!15}\line(1,0){92.0}}
\put(18.0,58.5){\makebox(0,0)[r]{0.70}}
\put(22.0,73.0){\color{black!15}\line(1,0){92.0}}
\put(18.0,73.0){\makebox(0,0)[r]{0.75}}
\put(22.0,15.0){\line(1,0){92.0}}
\put(22.0,15.0){\line(0,1){58.0}}
\put(26,79){\makebox(0,0)[l]{\rule{1.8mm}{1.8mm} court architecture \quad \raisebox{-0.3mm}{\framebox(1.8,1.8){}} mechanism scenario}}
\put(68.0,21.1){\makebox(0,0){\color{black}\rule{2.1mm}{2.1mm}}}
{\color{black!45}\qbezier[8](66.5,22.5)(66.0,24.4)(65.5,26.3)}
\put(64.6,26.3){\makebox(0,0)[r]{CUR*}}
\put(70.3,43.7){\makebox(0,0){\color{black}\rule{1.7mm}{1.7mm}}}
{\color{black!45}\qbezier[8](71.6,45.0)(75.5,50.3)(79.4,55.7)}
\put(80.3,55.7){\makebox(0,0)[l]{EMR}}
\put(100.0,44.0){\makebox(0,0){\color{black}\rule{1.7mm}{1.7mm}}}
\put(103.2,46.0){\makebox(0,0)[l]{DUAL}}
\put(87.1,44.6){\makebox(0,0){\color{black}\rule{1.7mm}{1.7mm}}}
{\color{black!45}\qbezier[8](85.8,45.8)(86.2,48.2)(86.6,50.6)}
\put(86.6,50.6){\makebox(0,0)[c]{COUNC}}
\put(71.7,40.2){\makebox(0,0){\color{black}\rule{1.7mm}{1.7mm}}}
{\color{black!45}\qbezier[8](72.9,39.0)(74.6,36.3)(76.2,33.6)}
\put(77.1,33.6){\makebox(0,0)[l]{OVR}}
\put(89.2,47.8){\makebox(0,0){\color{black}\rule{1.7mm}{1.7mm}}}
{\color{black!45}\qbezier[8](90.4,49.0)(91.3,50.4)(92.3,51.8)}
\put(93.2,51.8){\makebox(0,0)[l]{HYB}}
\put(46.2,60.3){\color{black}\framebox(1.7,1.7){}}
{\color{black!45}\qbezier[8](45.8,62.4)(44.7,63.8)(43.6,65.3)}
\put(42.7,65.3){\makebox(0,0)[r]{WFR}}
\put(51.7,51.9){\color{black}\framebox(1.7,1.7){}}
{\color{black!45}\qbezier[8](51.3,54.0)(50.1,53.8)(48.9,53.7)}
\put(48.0,53.7){\makebox(0,0)[r]{SUS}}
\put(64.2,48.1){\color{black}\framebox(1.7,1.7){}}
{\color{black!45}\qbezier[8](63.8,47.7)(60.8,47.4)(57.9,47.1)}
\put(57.0,47.1){\makebox(0,0)[r]{PRE}}
\put(63.9,47.5){\color{black}\framebox(1.7,1.7){}}
{\color{black!45}\qbezier[8](66.0,49.6)(67.4,50.4)(68.7,51.1)}
\put(69.6,51.1){\makebox(0,0)[l]{ABS}}
\put(82.3,37.9){\color{black}\framebox(1.7,1.7){}}
{\color{black!45}\qbezier[8](84.4,37.5)(85.4,35.7)(86.3,33.8)}
\put(87.2,33.8){\makebox(0,0)[l]{OMB}}
\put(82.8,42.0){\color{black}\framebox(1.7,1.7){}}
{\color{black!45}\qbezier[8](82.4,41.6)(81.5,39.7)(80.5,37.8)}
\put(79.6,37.8){\makebox(0,0)[r]{DEF}}
\put(64.2,55.0){\color{black}\framebox(1.7,1.7){}}
{\color{black!45}\qbezier[8](66.3,57.1)(67.2,61.3)(68.1,65.4)}
\put(69.0,65.4){\makebox(0,0)[l]{RIS}}
\put(54.3,51.9){\color{black}\framebox(1.7,1.7){}}
{\color{black!45}\qbezier[8](53.9,54.0)(54.2,57.8)(54.6,61.7)}
\put(54.6,61.7){\makebox(0,0)[c]{MLR}}
\put(68,4){\makebox(0,0){total institutional cost}}
\put(4,45){\rotatebox{90}{\makebox(0,0){democratic constitutionalism}}}
\end{picture}
\endgroup

\inlinefigurecaption{Baseline institutional cost versus democratic constitutionalism. The vertical axis is a synthetic diagnostic combining rights, responsiveness, legitimacy, compliance, implementation, access, response credibility, conflict, veto-relocation risk, feasibility, and culture sensitivity. CUR* marks the current-style scenario.}{fig:cost-score}
\end{inlinefigureblock}

Figure~\ref{fig:mechanism-tradeoff} focuses only on synthetic mechanisms. The desired quadrant is high democratic constitutionalism and low veto-relocation risk. This figure is not a ranking of legal systems; it is a stress test of whether a proposed device protects rights and preserves democratic response without hiding final policy control in a less accountable stage.

\begin{figure}[t]
\centering
% Auto-generated by paper/scripts/generate_figures.py
\begingroup
\setlength{\unitlength}{1mm}
\begin{picture}(128,86)
\scriptsize
\put(22.0,15.0){\color{black!15}\line(0,1){58.0}}
\put(22.0,11.0){\makebox(0,0){0.00}}
\put(47.1,15.0){\color{black!15}\line(0,1){58.0}}
\put(47.1,11.0){\makebox(0,0){0.15}}
\put(72.2,15.0){\color{black!15}\line(0,1){58.0}}
\put(72.2,11.0){\makebox(0,0){0.30}}
\put(97.3,15.0){\color{black!15}\line(0,1){58.0}}
\put(97.3,11.0){\makebox(0,0){0.45}}
\put(22.0,15.0){\color{black!15}\line(1,0){92.0}}
\put(18.0,15.0){\makebox(0,0)[r]{0.55}}
\put(22.0,29.5){\color{black!15}\line(1,0){92.0}}
\put(18.0,29.5){\makebox(0,0)[r]{0.60}}
\put(22.0,44.0){\color{black!15}\line(1,0){92.0}}
\put(18.0,44.0){\makebox(0,0)[r]{0.65}}
\put(22.0,58.5){\color{black!15}\line(1,0){92.0}}
\put(18.0,58.5){\makebox(0,0)[r]{0.70}}
\put(22.0,73.0){\color{black!15}\line(1,0){92.0}}
\put(18.0,73.0){\makebox(0,0)[r]{0.75}}
\put(22.0,15.0){\line(1,0){92.0}}
\put(22.0,15.0){\line(0,1){58.0}}
\put(51.6,61.1){\makebox(0,0){\color{black}\rule{1.8mm}{1.8mm}}}
{\color{black!45}\qbezier[8](50.4,62.3)(49.0,64.1)(47.6,65.9)}
\put(47.6,65.9){\makebox(0,0)[r]{WFR}}
\put(56.6,52.7){\makebox(0,0){\color{black}\rule{1.8mm}{1.8mm}}}
{\color{black!45}\qbezier[8](57.8,51.5)(59.2,49.9)(60.6,48.3)}
\put(60.6,48.3){\makebox(0,0)[l]{SUS}}
\put(81.4,45.5){\makebox(0,0){\color{black}\rule{1.8mm}{1.8mm}}}
{\color{black!45}\qbezier[8](82.6,46.7)(84.0,48.4)(85.4,50.1)}
\put(85.4,50.1){\makebox(0,0)[l]{OVC}}
\put(107.5,48.9){\makebox(0,0){\color{black}\rule{1.8mm}{1.8mm}}}
{\color{black!45}\qbezier[8](106.3,47.7)(104.9,45.9)(103.5,44.1)}
\put(103.5,44.1){\makebox(0,0)[r]{PRE}}
\put(74.2,48.4){\makebox(0,0){\color{black}\rule{1.8mm}{1.8mm}}}
{\color{black!45}\qbezier[8](75.4,49.6)(76.9,49.5)(78.4,49.4)}
\put(78.4,49.4){\makebox(0,0)[l]{ABS}}
\put(109.8,38.8){\makebox(0,0){\color{black}\rule{1.8mm}{1.8mm}}}
{\color{black!45}\qbezier[8](111.0,37.6)(112.5,35.8)(114.0,34.0)}
\put(114.0,34.0){\makebox(0,0)[l]{OMB}}
\put(108.0,42.8){\makebox(0,0){\color{black}\rule{1.8mm}{1.8mm}}}
{\color{black!45}\qbezier[8](106.8,41.6)(105.3,39.8)(103.8,38.0)}
\put(103.8,38.0){\makebox(0,0)[r]{DEF}}
\put(106.0,55.9){\makebox(0,0){\color{black}\rule{1.8mm}{1.8mm}}}
{\color{black!45}\qbezier[8](104.8,57.1)(103.3,58.9)(101.8,60.7)}
\put(101.8,60.7){\makebox(0,0)[r]{RIS}}
\put(59.5,52.7){\makebox(0,0){\color{black}\rule{1.8mm}{1.8mm}}}
{\color{black!45}\qbezier[8](60.7,53.9)(62.1,55.7)(63.5,57.5)}
\put(63.5,57.5){\makebox(0,0)[l]{MLR}}
\put(68,4){\makebox(0,0){veto-relocation risk}}
\put(4,45){\rotatebox{90}{\makebox(0,0){democratic constitutionalism}}}
\end{picture}
\endgroup

\caption{Synthetic mechanism tradeoff: democratic constitutionalism versus veto-relocation risk.}
\label{fig:mechanism-tradeoff}
\end{figure}

\section{Limitations and Next Steps}
\label{sec:limits}

The simulator is intentionally abstract. It compresses doctrinal variation, legal culture, party systems, federal design, and executive implementation into variables that can be compared across institutional scenarios. It does not yet model litigant strategy, judicial opinion assignment, actual precedent networks, budget appropriations, media framing, or reform backlash as independently strategic actors. Case-selection access, legislative-response credibility, legislative-response timing, government repeat-player advantage, implementation capacity, legal-transplant feasibility, and political-culture sensitivity are diagnostic constructs, not measured country-level estimates. Implementation cost is normalized to the simulator scale and should not be read as a budget forecast. The source-specific validation surface is narrow; most comparative rows remain stress-test assumptions until paired with reproducible empirical datasets.

The largest inferential risk is circularity between mechanism definitions and outcome labels. A rule requiring reasons will mechanically improve a transparency component, and a no-relief-without-merits rule will mechanically reduce a merits-displacing emergency component. Those results are useful for checking implementation consistency, but they are not the substantive payoff. The stronger evidence comes only when direct rule effects propagate into less mechanical diagnostics such as rights-claimant success, compliance, constitutional-conflict index, modeled public-legitimacy proxy, legislative response, delay, and cost. A second risk is parameter arbitrariness: the weights are explicit and sensitivity-tested, but they are not estimated coefficients. A third risk is comparative overreach. The validation checks do not identify real counterfactual effects for named courts, and several mechanisms are grounded more heavily in U.S. emergency-docket concerns than in fully coded comparative datasets. The paper therefore treats the current results as a structured thought experiment with plausibility checks, not as a validated prediction engine.

The most valuable empirical extensions are richer case-flow validation against specific courts and periods; denominator-backed evidence for emergency applications, merits follow-through, legislative response, compliance, enforcement, transplant feasibility, political culture, and case-selection access; and externally sourced indicators for party fragmentation, government control, electoral time pressure, civil-society capacity, public trust, court throughput, and administrative implementation capacity. Future empirical work should replace synthesized targets with denominator-backed datasets, model political-culture variables with external indicators rather than internal assumptions, add explicit implementation budgets for access institutions, and test whether legislative response cycles are credible under different party-system and compliance conditions.

\section*{Supplementary Material}
The supplementary replication package is designed to include the simulator source and tests, calibration and benchmark inputs, source-observation documentation, generated reports, figure- and table-generation scripts, and reproduction commands. Full case-level and interval exports are reproducible from those commands and should be included in the final public deposit or supplied separately if requested.

\section*{Funding Declaration}
No external funding is reported for this manuscript.

\section*{Competing Interests Declaration}
Competing interests: The author declares none.

\section*{Data Availability Statement}
A replication archive will be supplied with the submitted manuscript and, if accepted, deposited in an appropriate public repository. The archive will include source code, model inputs, calibration target-construction inputs, aggregate campaign reports, manuscript source, and commands needed to reproduce the reported outputs. Full case-level and interval outputs are reproducible from the same commands and should be included in the final public repository or supplied separately if requested during review.

\section*{Ethical Statement}
This project uses synthetic simulation outputs and public institutional benchmarks. It does not involve human-subjects data.

\section*{AI Tools Declaration}
OpenAI ChatGPT Deep Research and OpenAI Codex were used from May 1--9, 2026. ChatGPT Deep Research was accessed through the ChatGPT web interface at \url{https://chatgpt.com}; the product interface did not expose a more specific model version. It was used for source-discovery prompts, venue-format research, and research-question drafting. OpenAI Codex was accessed through the Codex desktop interface using GPT-5-family coding assistance as exposed by the local application. Codex was used for code editing, LaTeX drafting, deterministic figure and table generation, replication checks, and manuscript revision. The tools were not trained, fine-tuned, or otherwise modified for this project, and they were not listed as authors. AI outputs were not treated as sources; cited source claims, model choices, figures, code, and manuscript text remain subject to author review and revision.

\begin{thebibliography}{}
% Auto-generated by paper/scripts/generate_tables.py

\bibitem[Baude(2015)]{baude2015}
Baude, William. 2015. ``Foreword: The Supreme Court's Shadow Docket.'' \url{https://chicagounbound.uchicago.edu/journal_articles/8067/}. New York University Journal of Law and Liberty 9: 1.

\bibitem[Bickel(1962)]{bickel1962}
Bickel, Alexander M. 1962. \textit{The Least Dangerous Branch: The Supreme Court at the Bar of Politics}. New Haven, CT: Yale University Press.

\bibitem[Epp(1998)]{epp1998}
Epp, Charles R. 1998. \textit{The Rights Revolution: Lawyers, Activists, and Supreme Courts in Comparative Perspective}. Chicago: University of Chicago Press.

\bibitem[Epstein and Knight(1998)]{epsteinKnight1998}
Epstein, Lee, and Jack Knight. 1998. \textit{The Choices Justices Make}. Washington, DC: CQ Press.

\bibitem[Gallup(2024)]{gallupCourtApproval2024}
Gallup. 2024. ``Approval of U.S. Supreme Court Stalled Near Historical Low.'' \url{https://news.gallup.com/poll/647834/approval-supreme-court-stalled-near-historical-low.aspx}. Accessed May 1, 2026.

\bibitem[Gardbaum(2013)]{gardbaum2013}
Gardbaum, Stephen. 2013. \textit{The New Commonwealth Model of Constitutionalism: Theory and Practice}. Cambridge: Cambridge University Press.

\bibitem[Ginsburg(2003)]{ginsburg2003}
Ginsburg, Tom. 2003. \textit{Judicial Review in New Democracies: Constitutional Courts in Asian Cases}. Cambridge: Cambridge University Press.

\bibitem[Ginsburg and Versteeg(2014)]{ginsburgVersteeg2014}
Ginsburg, Tom, and Mila Versteeg. 2014. ``Why Do Countries Adopt Constitutional Review?.'' \url{https://academic.oup.com/jleo/article/30/3/587/881605}. Journal of Law, Economics, and Organization 30(3): 587--622.

\bibitem[Helmke(2005)]{helmke2005}
Helmke, Gretchen. 2005. \textit{Courts under Constraints: Judges, Generals, and Presidents in Argentina}. Cambridge: Cambridge University Press.

\bibitem[Hirschl(2004)]{hirschl2004}
Hirschl, Ran. 2004. \textit{Towards Juristocracy: The Origins and Consequences of the New Constitutionalism}. Cambridge, MA: Harvard University Press.

\bibitem[Kapiszewski et al.(2013)]{kapiszewski2013}
Kapiszewski, Diana, Gordon Silverstein, and Robert A. Kagan. 2013. \textit{Consequential Courts: Judicial Roles in Global Perspective}. Cambridge: Cambridge University Press.

\bibitem[SCOTUSblog(2025)]{scotusblogInterim2024}
SCOTUSblog. 2025. ``Interim Docket 2024.'' \url{https://www.scotusblog.com/cases/interim-docket/2024/}. Accessed May 1, 2026.

\bibitem[Staton(2010)]{staton2010}
Staton, Jeffrey K. 2010. \textit{Judicial Power and Strategic Communication in Mexico}. Cambridge: Cambridge University Press.

\bibitem[Stone Sweet(2000)]{stoneSweet2000}
Stone Sweet, Alec. 2000. \textit{Governing with Judges: Constitutional Politics in Europe}. Oxford: Oxford University Press.

\bibitem[Supreme Court Database(2025a)]{scdb2025}
Supreme Court Database. 2025a. ``Modern Database: 2025 Release 01.'' \url{https://scdb.la.psu.edu/data/2025-release-01/}. Accessed May 1, 2026.

\bibitem[Supreme Court Database(2025b)]{scdbIssueCodebook}
Supreme Court Database. 2025b. ``Online Codebook: Issue.'' \url{https://scdb.la.psu.edu/online-codebook/issue/}. Accessed May 1, 2026.

\bibitem[Tushnet(1999)]{tushnet1999}
Tushnet, Mark. 1999. \textit{Taking the Constitution Away from the Courts}. Princeton, NJ: Princeton University Press.

\bibitem[Vanberg(2005)]{vanberg2005}
Vanberg, Georg. 2005. \textit{The Politics of Constitutional Review in Germany}. Cambridge: Cambridge University Press.

\bibitem[Vladeck(2023)]{vladeck2023}
Vladeck, Stephen I. 2023. \textit{The Shadow Docket: How the Supreme Court Uses Stealth Rulings to Amass Power and Undermine the Republic}. New York: Basic Books.

\end{thebibliography}

\end{document}
